\documentclass[11pt]{amsart}
\input{glyphtounicode}
\pdfgentounicode=1

\usepackage[T1]{fontenc}
\usepackage[utf8]{inputenc}
\usepackage{amsmath,amssymb}
\usepackage{array,booktabs,tabularx}
\usepackage{graphicx}
\usepackage{geometry}
\usepackage{xcolor}
\usepackage{hyperref}
\usepackage{enumitem}
\geometry{margin=1in}
\setlength{\emergencystretch}{3em}
\sloppy
\hypersetup{
  colorlinks=true,
  linkcolor=blue!45!black,
  citecolor=blue!45!black,
  urlcolor=blue!45!black,
  pdftitle={Spectral A-N v18: Software Architecture and Reproducibility Companion},
  pdfauthor={Eliuth Chavero Jasso},
  pdfsubject={Software companion for the v18 certification suite},
  pdfkeywords={reproducibility, hashing, job scheduler, CPU/GPU parity, CI}
}

\title[Spectral A-N v18: Software and Reproducibility]{Spectral A-N v18: Software Architecture and Reproducibility Companion}
\author{Eliuth Chavero Jasso}
\address{Independent Researcher, Apizaco, Tlaxcala, Mexico}
\thanks{Companion to ``A Fail-Closed Certification Framework for Cyclic N-Ary Laws, Projectors, and Multiresolution Tensor Structure'' (\texttt{papers/a\_to\_n\_certification\_v18}). Documents interfaces and evidence contracts; makes no independent mathematical claim.}
\date{30 July 2026}

\begin{document}
\maketitle

\begin{abstract}
This companion documents the software architecture, artifact schemas,
hashing/provenance contracts, and reproducibility discipline of the v18
spectral certification suite (\texttt{spectral/certification\_v18/}, 85
passing tests as of this writing). It records a real, reproducible
CPU/GPU parity result on the confirmed hardware (RTX PRO 5000 Blackwell,
24\,GB VRAM), a concrete process-global TF32 contamination hazard found and
fixed, and the current gaps in continuous integration and packaging for
this new track.
\end{abstract}

\section{Repository layout}
\begin{center}
\small
\begin{tabular}{ll}
\toprule
Path & Contents \\
\midrule
\texttt{spectral/legacy/v17/} & Immutable, hash-verified legacy script + run-log copies \\
\texttt{spectral/certification\_v18/gates.py} & Typed-state enum, fail-closed combination logic \\
\texttt{spectral/certification\_v18/config.py} & Certification-mode contract, enforced via \texttt{validate()} \\
\texttt{spectral/certification\_v18/model.py} & Fresh CP-law/projector geometry (not imported from legacy) \\
\texttt{spectral/certification\_v18/blocks/} & One module + one \texttt{BLOCK\_*\_FINDINGS.md} per block \\
\texttt{spectral/certification\_v18/hardware/} & Certification-mode setup, resumable job queue \\
\texttt{spectral/certification\_v18/artifacts/} & Saved experiment result JSON \\
\texttt{spectral/certification\_v18/tests/} & 85 tests, \texttt{pytest spectral/certification\_v18/tests} \\
\bottomrule
\end{tabular}
\end{center}

This track is deliberately NOT reimported into \texttt{src/seion\_core}: it
is scope class \texttt{SPECTRAL\_LEGACY\_TRACK}
(\texttt{claims/scope\_registry\_v4.yaml}), evidentially separate from
\texttt{CANONICAL\_FINITE\_CORE} by design.

\section{Why \texttt{model.py} is a fresh reimplementation, not an import}
The legacy script sets \texttt{torch.backends.cuda.matmul.allow\_tf32 =
True} and \texttt{torch.set\_float32\_matmul\_precision("high")}
unconditionally at import time when CUDA is available. These are
process-global flags, not scoped to the importing module. A certification-mode
process that imported the legacy script for convenience would have its
TF32 discipline silently overridden by an unrelated import. `model.py`
reimplements the needed primitives (CP product, projector, associator,
reduced curvature) independently, with its own tests
(\texttt{test\_model.py}) checked against the legacy formulas by direct
line-number citation, and \texttt{hardware/certification\_mode.py}
provides \texttt{enter\_certification\_mode()} / \\
\texttt{assert\_certification\_discipline()} as the single, idempotent,
re-assertable enforcement point.

\section{Typed-gate schema}
\texttt{gates.py} defines \texttt{TypedStatus} (10 states) and
\texttt{CRITICAL\_GATES} (8 gates, each a tuple of contributing block
names). \texttt{combine\_gate\_status} takes the minimum over a gate's
blocks (excluding \texttt{NOT\_APPLICABLE}, which is tracked separately
and never silently counted as passing). \texttt{assign\_block\_status}
is the single enforcement point preventing a screening-mode run from
being assigned a certificate-tier status: it raises
\texttt{ScreeningCertificateViolation} rather than downgrading silently.

\section{CPU/GPU parity: a real, executed result}
On the confirmed hardware (Intel Core Ultra 9 285HX, 24 logical cores,
128\,GB RAM; NVIDIA RTX PRO 5000 Blackwell Laptop GPU, 24\,GB VRAM, CUDA
12.8, driver 573.49), we constructed a \texttt{SpectralModelV18} instance
on CPU (float64, seed 0), transferred its exact parameter/buffer values
to an independently-constructed CUDA instance via \texttt{state\_dict}
(not \texttt{.to(device)}, which does not update the plain Python
\texttt{device}/\texttt{dtype} string attributes stored on
\texttt{CyclicCPProduct} --- a second real hazard found during this
exercise, documented in \texttt{model.py}), and computed Block A/B
quantities on both devices with TF32 explicitly disabled and
deterministic algorithms enabled. Results:
\begin{center}
\small
\begin{tabular}{lrr}
\toprule
Quantity & CPU & CUDA \\
\midrule
idempotence residual & $1.312\times10^{-15}$ & $1.198\times10^{-15}$ \\
self-adjointness residual & $0.0$ (exact) & $2.220\times10^{-16}$ \\
comm.\ unexplained-rel (random init) & $1.0000380759060201$ & $1.0000380759060197$ \\
wall time & $19.1$\,ms & $835.1$\,ms \\
\bottomrule
\end{tabular}
\end{center}
The unexplained-rel figure (the only quantity here large enough for a
relative-difference comparison to be meaningful) agrees to
$4.4\times10^{-16}$ relative difference --- CPU/GPU parity confirmed at
float64 for this computation. The wall-time comparison is an honest
negative result for accelerating problems at this scale: GPU kernel-launch
overhead dominates for $n=24$, and CPU is $\sim 44\times$ faster in this
single case. Real GPU throughput advantage should be expected only once
batch sizes or ambient dimensions are large enough to amortize that
overhead --- not exercised in this pass; the sweep/scaling phase
(Section~\ref{sec:gaps}) is where that would be measured.

\section{Resumable job queue}
\texttt{hardware/job\_queue.py} implements an append-only, JSONL-backed
ledger keyed by \texttt{scientific\_instance\_id} (the mathematical
configuration) versus \texttt{execution\_id} (one specific attempt),
recording seed, precision, hardware, a config hash, a script hash,
checkpoint lineage (parent \texttt{execution\_id}), status, retry count,
and output hashes. History is never mutated, only appended to, matching
the pattern already established for the separate \texttt{CANONICAL\_FINITE\_CORE}
track's \texttt{src/seion\_core/governance/runs.py} (different schema,
same non-destructive philosophy, since this track's artifact set does not
match that track's \texttt{REQUIRED\_RUN\_ARTIFACTS} contract). Tested:
submit/complete, retry-increments-count while preserving prior history,
and checkpoint-lineage traversal.

\section{Legacy provenance}
\texttt{spectral/certification\_v18/tools/ingest\_legacy.py} hashes both
named legacy files (sha256, verified after copy), hashes all 18 on-disk
run directories, parses the 9-run text log, cross-references it against
on-disk \texttt{summary.json} files, reconstructs resume lineage, and
reclassifies every historical run against the v18 typed-gate vocabulary.
Output: \texttt{spectral/legacy/v17/legacy\_a\_to\_n\_manifest.yaml},
\texttt{legacy\_run\_lineage.json}, \texttt{legacy\_run\_dedup\_report.md},
\texttt{legacy\_claim\_reclassification.yaml}.

\section{Gaps and follow-up (honestly listed, not hidden)}
\label{sec:gaps}
\begin{itemize}[leftmargin=*]
\item No CI workflow for \texttt{spectral/certification\_v18/} exists yet
  in \texttt{.github/}; the 85 tests are currently run manually only.
\item No SBOM or packaging metadata specific to this track has been
  generated; the track depends on the same \texttt{pyproject.toml}
  dependencies (\texttt{torch}, \texttt{numpy}) as the rest of the
  repository, declared under the existing \texttt{research} extra.
\item GPU throughput scaling (large batch/dimension sweeps) was not
  measured --- only correctness parity at small scale.
\item Clean-environment reproduction (a fresh virtualenv/container
  install-and-rerun) was not exercised this pass; all runs used the
  already-configured local environment.
\end{itemize}

\end{document}
